\chapter{Lernziele Kryptologie}

Die folgenden Lernziele geben Ihnen einen Überblick über die zentralen Konzepte und Fähigkeiten, die in den einzelnen Kapiteln dieses Skripts behandelt werden.

\section*{Kapitel 1: Geheimschriften und einfache Kryptosysteme}
\begin{todolist}
  \item Ich kenne die Grundbegriffe der Kryptologie: Klartext, Kryptotext, Schlüssel, Verschlüsselung, Entschlüsselung, Kryptosystem sowie die Abgrenzung zwischen Kryptographie, Kryptoanalyse und Kryptologie.
  \item Ich kenne die vier zentralen Schutzziele: Vertraulichkeit, Integrität, Authentizität und Verbindlichkeit (Nichtabstreitbarkeit) und kann sie an Beispielen erläutern.
  \item Ich kann das Kerckhoffs-Prinzip erklären und begründen, warum Sicherheit ausschliesslich auf der Geheimhaltung des Schlüssels und nicht des Algorithmus beruhen darf.
  \item Ich kann einfache historische Verfahren (Skytale, Caesar, monoalphabetische Substitution) zur Ver- und Entschlüsselung anwenden.
  \item Ich kann einen Caesar-Kryptotext durch Ausprobieren aller 25 möglichen Verschiebungen (Brute-Force) brechen.
  \item Ich kann das Prinzip der Häufigkeitsanalyse natürlicher Sprachen erklären und anwenden, um monoalphabetisch verschlüsselte Texte zu knacken.
\end{todolist}

\section*{Kapitel 2: Vigenère und One-Time-Pad}
\begin{todolist}
  \item Ich kann den Unterschied zwischen monoalphabetischer und polyalphabetischer Substitution erklären.
  \item Ich kann Nachrichten mit der Vigenère-Chiffre bei gegebenem Schlüsselwort ver- und entschlüsseln.
  \item Ich kann erklären, wie die Schlüssellänge einer Vigenère-Chiffre bestimmt werden kann (Kasiski-Test und Friedmansche Charakteristik / Koinzidenzindex).
  \item Ich kann erläutern, wie Vigenère bei bekannter Schlüssellänge durch Zerlegung in Teiltexte und gruppenweise Häufigkeitsanalyse gebrochen werden kann.
  \item Ich kann begründen, warum das One-Time-Pad (OTP) perfekte Sicherheit nach Shannon bietet (echter Zufall, Schlüssellänge mindestens Textlänge, Einmaligkeit).
  \item Ich kann erklären, warum die Wiederverwendung eines OTP-Schlüssels (Two-Time-Pad) zu Informationsleckage führt.
  \item Ich kenne die Eigenschaften der XOR-Operation ($\oplus$: Kommutativität, Assoziativität, neutrales Element 0, Selbstinversität: $a \oplus a = 0$) und kann das binäre OTP für Ver- und Entschlüsselung anwenden.
\end{todolist}

\section*{Kapitel 3: Schlüsseltausch-Verfahren}
\begin{todolist}
  \item Ich kann das Problem des Schlüsselaustauschs bei symmetrischen Verfahren erläutern.
  \item Ich kann das Shamir Three-Pass Protocol (Drei-Wege-Schlüsseltausch) schrittweise durchführen und erklären, warum es bei ausschliesslicher Verwendung von XOR unsicher ist ($k_A \oplus k_AB \oplus k_B = t$).
  \item Ich kann die Bedeutung des Definitionszeichens (:=) sowie des Kongruenzzeichens ($\equiv$) erklären.
  \item Ich kann das Diffie-Hellman-Merkle-Verfahren (DHM) mit kleinen Zahlen durchführen und den gemeinsamen Schlüssel berechnen ($K = g^{ab} \bmod p$).
  \item Ich kann den Man-in-the-Middle-Angriff (MitM) auf unauthentifizierte Schlüsseltauschverfahren beschreiben.
\end{todolist}

\section*{Kapitel 4: Asymmetrische Kryptosysteme und Digitale Signaturen}
\begin{todolist}
  \item Ich kann das Prinzip der Public-Key-Kryptographie (Schlüsselpaar aus öffentlichem und privatem Schlüssel) und den Unterschied zu symmetrischen Verfahren erklären.
  \item Ich kann die Schritte des RSA-Verfahrens (Schlüsselerzeugung mit Primzahlen $p, q$, Verschlüsselung $c = m^e \bmod n$, Entschlüsselung $m = c^d \bmod n$) mit kleinen Zahlen durchführen.
  \item Ich kann die Sicherheitsgrundlage von RSA (Schwierigkeit der Primfaktorzerlegung von $n = p \cdot q$) erläutern.
  \item Ich kann das Prinzip digitaler Signaturen mit RSA durch Rollenumkehr der Schlüssel erklären (Signieren mit privatem Schlüssel: $s = m^d \bmod n$, Verifizieren mit öffentlichem Schlüssel: $s^e \equiv m \pmod n$).
  \item Ich kann erklären, welche Schutzziele digitale Signaturen gewährleisten (Authentizität der Absenderin und Integrität der Nachricht).
  % \item Ich kann die Rolle kryptographischer Hashfunktionen (digitaler Fingerabdruck, Lawineneffekt, Kollisionsresistenz) erklären und begründen, warum in der Praxis Hashwerte signiert werden.
  \item Ich kann den Unterschied zwischen einer reinen digitalen Signatur (Klartextübertragung) und kombinierter Signatur und Verschlüsselung (\emph{Sign-then-Encrypt}) anhand ihrer Schutzziele erklären und anwenden.
\end{todolist}

\section*{Kapitel 5: Digitale Zertifikate und Public-Key-Infrastrukturen (PKI)}
\begin{todolist}
  \item Ich kann das Vertrauensproblem unbeglaubigter öffentlicher Schlüssel in der Praxis beschreiben (fehlende Bindung zwischen Schlüssel und Person/Domain).
  \item Ich kann erklären, wie ein digitales Zertifikat und eine Zertifizierungsstelle (CA) dieses Problem lösen, indem die CA die Bindung eines Schlüssels an eine Identität signiert.
  \item Ich kann die zwei verschiedenen Signaturen in einer PKI unterscheiden: Signatur der CA auf dem Zertifikat vs. Signatur des Besitzers auf einer Nachricht.
  \item Ich kann erklären, warum eine erfolgreiche Zertifikats- und Signaturprüfung erst zusammen mit dem Vergleich der zertifizierten Identität (Name bzw.\ Domain) belegt, dass die erwartete Person signiert hat.
  \item Ich kann das Missverständnis aufklären, dass ein gültiges Zertifikat nicht die allgemeine Vertrauenswürdigkeit eines Anbieters beweist, sondern ausschliesslich die Echtheit der Domainbindung (Phishing-Gefahr).
  \item Ich kann den Aufbau einer hierarchischen Zertifikatkette (Root-CA im Trust Store des Betriebssystems/Browsers, Intermediate-CA, Serverzertifikat) und den Sicherheitsvorteil der Kaskade erklären.
\end{todolist}
